\section{Related Work}
\label{sec:related}

\para{Number representations in LLMs.}
A growing body of work investigates how language models encode numerical information.
\citet{levy2024language} demonstrate that models use per-digit circular representations in base 10, achieving 91--96\% probe accuracy for individual digits in Llama~3 and \mistral.
\citet{kadlcik2025pretrained} introduce sinusoidal-basis probes that achieve near-perfect accuracy in decoding numeric values, overturning earlier findings of imprecise number representations.
\citet{stefanik2025unravelling} extend this work by showing that sinusoidal representations are universal across models, layers, and contexts, with cross-model Fourier frequency overlap of 1.0 for the top frequencies.
\citet{davies2025language} caution that embedding-level representations are far from continuous scalar encodings, with string artifacts and tokenization effects dominating the embedding space.
\citet{marjieh2025number} find that LLM number representations are entangled, blending string-level properties (edit distance) with numerical properties (log-linear magnitude).
Our work extends this line of research to French vigesimal numbers, testing whether the base-10 internal representations documented in prior work accommodate numbers whose linguistic form encodes base-20 arithmetic.

\para{Multilingual numerical reasoning.}
\citet{bhattacharya2025investigating} study cross-linguistic numeral puzzles and find that LLMs fail not because of mathematical difficulty but because they cannot infer implicit compositional operators (addition, multiplication) from numeral structure.
Their work explicitly identifies French as an example vigesimal system, though they do not probe French representations directly.
\citet{rahman2025fragile} demonstrate that LLM numerical reasoning is fragile---models succeed at deterministic tasks but fail at combinatorial puzzles, revealing pattern matching over genuine number sense.
Unlike these studies, which focus on behavioral evaluation, we complement behavioral analysis with representational probing to understand \emph{why} models succeed or fail.

\para{Tokenization and arithmetic.}
\citet{singh2024tokenization} show that tokenization direction significantly impacts arithmetic accuracy, with right-to-left tokenization improving performance by up to 20\%.
\citet{graphpku2024number} establish the NUPA benchmark with systematic evaluation of number representations and find that performance degrades with number length and that single-digit tokenizers outperform multi-digit ones.
\citet{counting2026mechanistic} analyze the mechanistic basis of counting in LLMs, finding logarithmic compression and specific attention heads that mediate count information.
These findings are directly relevant to French vigesimal numbers, which require 40\% more tokens than their decimal counterparts, creating additional computational overhead for the transformer.

\para{Positioning of our work.}
Prior work on number representations has focused almost exclusively on English and digit tokens.
We contribute the first study of French vigesimal representations, using Belgian French as a controlled comparison that isolates the effect of linguistic irregularity.
Our combined behavioral-representational approach reveals structure invisible to behavioral benchmarks alone.
